\section*{ตัวอย่างการเขียน Latex}
\subsection*{การใช้ตัวอักษรรูปแบบต่างๆ}

ใน LaTeX, การใช้คำสั่ง \textbf{ตัวหนา} \textit{ตัวเอียง} \underline{ข้อความที่มีเส้นใต้} และ \textbf{\textit{ตัวหนาและตัวเอียง}} สามารถใช้ได้กับข้อความภาษาไทยได้เหมือนกับภาษาอังกฤษ

การใช้เครื่องหมายคำพูด เช่น ``ข้อความที่ต้องการอ้างถึง''

\subsection*{การสร้างรายการแบบ description}
\begin{description}
\item[\textbf{First}:] \dotfill
\item[\textbf{Second}:] \dotfill
\item[\textbf{Third}:] \dotfill
\end{description}

\subsection*{การสร้างรายการแบบ enumerate}
\begin{enumerate}
\item \textbf{First}: \dotfill
\item \textbf{Second}: \dotfill
\item \textbf{Third}: \dotfill
\end{enumerate}

\subsection*{การสร้างรายการแบบ itemize}
\begin{itemize}
\item \textbf{First}: \dotfill
\item \textbf{Second}: \dotfill
\item \textbf{Third}: \dotfill
\end{itemize}

\subsection*{การอ้างอิง}
\dotfill 
จากงานวิจัยของ \cite{musikawan2022enhanced,kongsorot2023intrusion} แสดงให้เห็นว่า \dotfill

\cite{musikawan2022enhanced} แสดงให้เห็นว่า \dotfill \cite{kongsorot2023intrusion} แสดงให้เห็นว่า \dotfill

\subsection*{การแทรกสมการและการอ้างอิงในเนื้อหา}
\begin{equation}
\frac{dy}{dx}+P(x)y(x)=Q(x)y^n\quad\text{เมื่อ }n\neq 0
\label{eq:equation_1}
\end{equation}

\begin{equation}
\begin{pmatrix}
\hat{r}\\ \hat{\theta}\\ \hat{\varphi}
\end{pmatrix}
=
\begin{pmatrix}
\sin\theta\cos\varphi&\sin\theta\sin\varphi&\cos\theta\\
\cos\theta\cos\varphi&\cos\theta\sin\varphi&-\sin\theta\\
-\sin\varphi&\cos\varphi&0\\
\end{pmatrix}
\begin{pmatrix}
\hat{x}\\ \hat{y}\\ \hat{z}
\end{pmatrix}
\label{eq:equation_2}
\end{equation}

\begin{align}
\sum{F}&=ma\nonumber\\
-kx-bv&=m\frac{d^2x}{dt^2}\nonumber\\
\frac{d^2x}{dt^2}+\frac{b}{m}\dfrac{dx}{dt}+\frac{k}{m}x&=0
\label{eq:equation_3}
\end{align}

จากสมการที่ \ref{eq:equation_1} แสดงให้เห็นว่า \dotfill

จากสมการที่ \ref{eq:equation_1} และ \ref{eq:equation_2} แสดงให้เห็นว่า \dotfill

จากสมการที่ \ref{eq:equation_1} \ref{eq:equation_2} และ \ref{eq:equation_3} แสดงให้เห็นว่า \dotfill

\subsection*{การแทรกรูปภาพและการอ้างอิงในเนื้อหา}
\begin{figure}[H]
\centering
\includegraphics[width=0.5\linewidth]{LaTeX.png}
\caption{ตัวอย่างการแทรกรูปภาพ}
\label{fig:my_figure}
\end{figure}

จากรูปที่ \ref{fig:my_figure} แสดงให้เห็นว่า \dotfill 

\subsection*{การแทรกตารางและการอ้างอิงในเนื้อหา}

\begin{table}[H]
\centering
\caption{ตัวอย่างการแทรกตาราง}
\label{tab:my_label_1}
\begin{tabular}{cccc}
\hline
คอลัมน์ที่ 0 & คอลัมน์ที่ 1 & คอลัมน์ที่ 2  & คอลัมน์ที่ 3 \\ \hline
- & - & - & - \\
- & - & - & - \\
- & - & - & - \\
- & - & - & - \\
\hline
\end{tabular}
\end{table}

\begin{table}[H]
\centering
\caption{ตัวอย่างการแทรกตาราง}
\label{tab:my_label_2}
\begin{tabular}{|c|c|}
\hline
\textbf{สาร} & \textbf{ความหนาแน่น (kg/m)}$^\mathbf{3}$\\ \hline
น้ำ & $1.0\times10^3$ \\ \hline
อลูมิเนียม & $2.7\times10^3$ \\ \hline
\end{tabular}
\end{table}

จากตารางที่ \ref{tab:my_label_1} แสดงให้เห็นว่า \dotfill 

จากตารางที่ \ref{tab:my_label_1} และ \ref{tab:my_label_2} แสดงให้เห็นว่า \dotfill 

\subsection*{การใช้ footnote}
การใช้ หมายเหตุเชิงเทียน (footnote) ในเอกสารมีข้อดีหลายประการ โดยเฉพาะในกรณีที่คุณต้องการแทรกข้อมูลเพิ่มเติมหรืออ้างอิงแหล่งที่มา (เช่น เว็บไซต์) โดยไม่ทำให้เนื้อหาหลักของเอกสารดูยุ่งเหยิงหรือยาวเกินไป

ตัวอย่าง

For more information on PyTorch, visit the official website\footnote{\url{https://pytorch.org}}.